\documentclass[../../main.tex]{subfiles}
\graphicspath{{../../image/}{../../figure/}}
\begin{document}

\section{Lax 激波容许性准则；压缩、超压缩与亚压缩激波}\label{section:ch08-lax容许性准则}

本节的目的是引入若干条件，使得激波与解中周围“较光滑”部分的相互作用是稳定的。稳定性的一种表现是：与激波相撞的小振幅光滑波，被吸收、透射和/或反射为小振幅波，而不会通过显著改变激波本身的强度或传播速度来影响其完整性。基于此，考虑上半平面上\eqref{align:ch08-守恒律一维} 的一个解：左侧常状态 $U_-$ 经激波 $x=st$ 与右侧常状态 $U_+$ 相连。特别地，Rankine--Hugoniot 跳跃条件\eqref{align:ch08-rh} 成立。设激波速度 $s$ 满足
\begin{align}
\begin{cases}
\lambda_1(U_-)\leqslant\cdots\leqslant\lambda_{i-1}(U_-)<s<\lambda_i(U_-)\leqslant\cdots\leqslant\lambda_n(U_-)\\
\lambda_1(U_+)\leqslant\cdots\leqslant\lambda_j(U_+)<s<\lambda_{j+1}(U_+)\leqslant\cdots\leqslant\lambda_n(U_+),
\end{cases}
\label{align:ch08-lax谱条件}
\end{align}
其中 $i=1,\cdots,n$，$j=1,\cdots,n$，并约定若 $i=1$ 则 $\lambda_{i-1}(U_-)=-\infty$，若 $j=n$ 则 $\lambda_{j+1}(U_+)=+\infty$。我们的方案是：通过对常状态 $U_-$ 与 $U_+$ 作轻微扰动来构造一族解。带有小振荡的\eqref{align:ch08-守恒律一维} 经典解，由线性化系统的解紧密逼近。因此，为在不陷入冗长技术细节的情况下达到预期结论，我们作简化假设：$\mathrm{D}F(U)$ 在 $|U-U_-|<\varepsilon$ 上为常数、等于 $\mathrm{D}F(U_-)$，并在 $|U-U_+|<\varepsilon$ 上为常数、等于 $\mathrm{D}F(U_+)$。于是在其上半平面寻找形如
\begin{align}
U(x,t)=
\begin{cases}
U_-+\sum_{k=1}^{n}\omega_k^-(x,t)R_k(U_-),&x<st+\sigma(t)\\
U_++\sum_{k=1}^{n}\omega_k^+(x,t)R_k(U_+),&x>st+\sigma(t),
\end{cases}
\label{align:ch08-扰动解}
\end{align}
的解，其中 $\omega_k^\pm(x,t)$ 与 $\sigma(t)$ 是 $C^1$ 函数，满足对某个 $a\ll1$ 有 $|\omega_k^\pm|<a\varepsilon$ 且 $|\dot\sigma|<a\varepsilon$。由于在 $U_\pm$ 邻域内 $F(U)=F(U_\pm)+\mathrm{D}F(U_\pm)(U-U_\pm)$，且 $R_k(U_\pm)$ 是 $\mathrm{D}F(U_\pm)$ 相应于特征值 $\lambda_k(U_\pm)$ 的特征向量，故\eqref{align:ch08-扰动解} 给出的 $U$ 对 $x\neq st+\sigma(t)$ 满足\eqref{align:ch08-守恒律一维}，只要对 $k=1,\cdots,n$ 有
\begin{align}
\begin{cases}
\partial_t\omega_k^-(x,t)+\lambda_k(U_-)\partial_x\omega_k^-(x,t)=0,&x<st+\sigma(t)\\
\partial_t\omega_k^+(x,t)+\lambda_k(U_+)\partial_x\omega_k^+(x,t)=0,&x>st+\sigma(t),
\end{cases}
\label{align:ch08-扰动方程}
\end{align}
并且， across 扰动激波 $x=st+\sigma(t)$ 它还满足 Rankine--Hugoniot 跳跃条件，当且仅当
\begin{align}
\sum_{k=1}^{n}[\lambda_k(U_+)-s-\dot\sigma]\omega_k^+R_k(U_+)
-\sum_{k=1}^{n}[\lambda_k(U_-)-s-\dot\sigma]\omega_k^-R_k(U_-)
=\dot\sigma(U_+-U_-).
\label{align:ch08-扰动rh}
\end{align}

由\eqref{align:ch08-扰动方程}，$\omega_k^\pm$ 沿 $k$-特征线为常数，这些特征线是斜率为 $\lambda_k(U_\pm)$ 的直线。特别地，由\eqref{align:ch08-lax谱条件}，$\omega_i^-,\cdots,\omega_n^-$ 与 $\omega_1^+,\cdots,\omega_j^+$ 在各处由其在 $t=0$ 处给定的初值完全确定。与之相反，$\omega_1^-,\cdots,\omega_{i-1}^-$ 与 $\omega_{j+1}^+,\cdots,\omega_n^+$ 部分由其在 $t=0$ 处给定的初值、部分由其沿激波上的值确定。后者受 Rankine--Hugoniot 条件\eqref{align:ch08-扰动rh} 约束，应把它看作关于 $n+i-j$ 个未知量 $\omega_1^-,\cdots,\omega_{i-1}^-,\omega_{j+1}^+,\cdots,\omega_n^+$ 及 $\dot\sigma$ 的 $n$ 个方程的方程组。

通常，激波以 $i=j$ 满足\eqref{align:ch08-lax谱条件}，于是
\begin{align}
\begin{cases}
\lambda_{i-1}(U_-)<s<\lambda_i(U_-)\\
\lambda_i(U_+)<s<\lambda_{i+1}(U_+).
\end{cases}
\label{align:ch08-压缩}
\end{align}
此时激波称为 \emph{压缩的}（compressive），该术语借自气体动力学。对压缩激波，\eqref{align:ch08-扰动rh} 中未知量个数为 $n$，与方程个数相匹配。若还满足
\begin{align}
\det\big[R_1(U_-),\cdots,R_{i-1}(U_-),U_+-U_-,R_{i+1}(U_+),\cdots,R_n(U_+)\big]\neq0,
\label{align:ch08-非退化}
\end{align}
则由隐函数定理容易证明：\eqref{align:ch08-扰动rh} 唯一确定 $\omega_1^-,\cdots,\omega_{i-1}^-,\omega_{j+1}^+,\cdots,\omega_n^+$ 及 $\dot\sigma$。在此条件下，一旦在允许范围内给定初值 $\omega_k^\pm(x,0)$，形如\eqref{align:ch08-扰动解} 的解 $U$ 便被唯一确定，该激波称为 \emph{演化的}（evolutionary）。

以 $i<j$ 满足\eqref{align:ch08-lax谱条件} 的激波称为 \emph{超压缩的}（overcompressive）。此时\eqref{align:ch08-扰动rh} 中方程个数超过未知量个数，因而形如\eqref{align:ch08-扰动解} 的解 $U$ 的存在性没有保证。尽管如此，此类激波确实在某些应用中出现。一个简单例子是系统\eqref{align:ch07-脐点模型}：回顾\hyperref[section:ch08-hugoniot轨迹]{8.2 节}中描述的其 Hugoniot 轨迹的形状，考虑一个速度为 $s$ 的激波，它在左侧连接位于单位圆上的状态 $(u_-,v_-)$，在右侧连接状态 $(u_+,v_+)=a(u_-,v_-)$，其中 $a$ 为某常数。由\eqref{align:ch07-脐点模型特征速度}，$\lambda_1(u_-,v_-)=1$，$\lambda_2(u_-,v_-)=3$，$\lambda_1(u_+,v_+)=a^2$，$\lambda_2(u_+,v_+)=3a^2$。进而 Rankine--Hugoniot 条件\eqref{align:ch08-umiblic} 给出 $s=a^2+a+1$。因此，若 $a\in(-\frac12,0)$，则 $\lambda_2(u_-,v_-)>\lambda_1(u_-,v_-)>s>\lambda_2(u_+,v_+)>\lambda_1(u_+,v_+)$，即该激波是超压缩的。

在相反情形，即\eqref{align:ch08-lax谱条件} 以 $i>j$ 成立时，激波称为 \emph{亚压缩的}（undercompressive）。

在该情形下，\eqref{align:ch08-扰动rh} 是欠定的，一般容许多个解。然而，当此类激波在应用中出现时，Rankine--Hugoniot 跳跃条件通常由形如 $G(U_-,U_+,s)=0$ 的方程（称为\emph{动力学关系}，kinetic relations）补充，从而使解唯一。

上述内容将在\hyperref[section:ch08-黏性激波剖面]{8.6 节}中进一步阐明。

条件\eqref{align:ch08-lax谱条件} 排除了以特征速度传播的激波，但此类激波确实存在。特别地，任一在左侧连接 $U_-$、在右侧连接 $U_+$、以速度 $s$ 传播的激波：若 $s=\lambda_i(U_-)$，称为左 $i$-接触间断；若 $s=\lambda_i(U_+)$，称为右 $i$-接触间断；若 $s=\lambda_i(U_-)=\lambda_i(U_+)$，简称为 $i$-接触间断。由\cref{theorem:ch08-线性退化激波曲线}，关联于线性退化特征族的任一弱激波必然是接触间断。

下文假设 $\lambda_i$ 是 $\mathrm{D}F$ 的单特征值，并把注意力集中于中等强度的压缩 $i$-激波。此时 $|s-\lambda_i|$ 与 $\lambda_{i+1}-\lambda_i$、$\lambda_i-\lambda_{i-1}$ 相比都很小，故\eqref{align:ch08-压缩} 中第一、第四个不等式恒成立。剩余两个不等式略加放宽以容许左和/或右 $i$-接触间断，合并为
\begin{align}
\lambda_i(U_+)\leqslant s\leqslant\lambda_i(U_-),
\label{align:ch08-lax-e}
\end{align}
这就是著名的 Lax $E$-条件。进而，由于 $U_+-U_-$ 与 $R_i(U_\pm)$ 近似共线，\eqref{align:ch08-非退化} 由双曲性蕴涵。

为说明起见，我们看几个例子，从单个守恒律\eqref{align:ch07-标量守恒律} 开始。特征速度为 $\lambda(u)=f'(u)$，故\eqref{align:ch08-lax-e} 取形式
\begin{align}
f'(u_+)\leqslant s\leqslant f'(u_-),
\label{align:ch08-标量lax}
\end{align}
其中 $s$ 是由 Rankine--Hugoniot 跳跃条件计算的激波速度：
\begin{align}
s=\frac{f(u_+)-f(u_-)}{u_+-u_-}.
\label{align:ch08-标量s}
\end{align}
读者注意到 $f'(u_-)$ 与 $f'(u_+)$ 是 $f$ 的图像在点 $(u_-,f(u_-))$ 与 $(u_+,f(u_+))$ 处的斜率，而 $s$ 是连接这两点的弦的斜率，便可立即领会\eqref{align:ch08-标量lax} 的几何解释。特别地，当\eqref{align:ch07-标量守恒律} 真非线性（即 $f''(u)\neq0$ 对所有 $u$）时，\eqref{align:ch08-标量lax} 化为：若 $f''(u)<0$ 则 $u_-<u_+$；若 $f''(u)>0$ 则 $u_->u_+$。

接下来考虑等熵弹性系统\eqref{align:ch07-p系统}。特征速度记录于\eqref{align:ch07-p系统特征值}，激波速度记录于\eqref{align:ch08-弹性s}，于是\eqref{align:ch08-lax-e} 对 1-激波或 2-激波分别取形式
\begin{align}
\sigma'(u_+)\geqslant\frac{\sigma(u_+)-\sigma(u_-)}{u_+-u_-}\geqslant\sigma'(u_-)
\quad\text{或}\quad
\sigma'(u_+)\leqslant\frac{\sigma(u_+)-\sigma(u_-)}{u_+-u_-}\leqslant\sigma'(u_-).
\label{align:ch08-弹性lax}
\end{align}
\eqref{align:ch08-弹性lax} 的几何解释同样清楚。当\eqref{align:ch07-p系统} 真非线性（即 $\sigma''(u)\neq0$ 对所有 $u$）时，\eqref{align:ch08-弹性lax} 化为：若 $\sigma''(u)>0$ 则 $u_-<u_+$ 或 $u_->u_+$；若 $\sigma''(u)<0$ 则 $u_->u_+$ 或 $u_-<u_+$。等价地，由\eqref{align:ch08-弹性rh}，以速度表示为：对两个激波族，若 $\sigma''(u)>0$ 则 $v_-<v_+$；若 $\sigma''(u)<0$ 则 $v_->v_+$。

类似分析适用于理想气体等熵流系统\eqref{align:ch07-理想气体}，其特征速度由\eqref{align:ch07-理想气体特征速度} 给出，并得到：一个在左侧连接状态 $(\rho_-,v_-)$、在右侧连接状态 $(\rho_+,v_+)$ 的 1-激波（或 2-激波）满足 Lax $E$-条件，当且仅当 $\rho_-<\rho_+$（或 $\rho_->\rho_+$）。换言之，容许激波阵面的通过压缩气体。

现在回到一般系统\eqref{align:ch08-守恒律一维}，考虑左侧状态 $U_-$ 与右侧状态 $U_+$ 由一个中等强度、速度为 $s$ 的 $i$-激波相连。设 $\lambda_i(U_-)$ 是 $\mathrm{D}F(U_-)$ 的单特征值，令 $W_i$ 表示过 $U_-$ 的 $i$-激波曲线（参见\cref{theorem:ch08-激波曲线}），于是 $U_-=W_i(0)$、$U_+=W_i(\tau)$。进而 $\lambda_i(U_-)=s_i(0)$、$s=s_i(\tau)$。我们证明：若 $\tau<0$ 且 $\dot s_i(\cdot)\geqslant0$ 在 $(\tau,0)$ 上成立，则该激波满足 Lax $E$-条件。确实，$\dot s_i(\cdot)\geqslant0$ 蕴涵 $s=s_i(\tau)\leqslant s_i(0)=\lambda_i(U_-)$，即\eqref{align:ch08-lax-e} 的右半部分。同时，只要\eqref{align:ch08-sdotpos} 与\eqref{align:ch08-sdotzero} 在 $\tau$ 处成立，$\dot s_i(\cdot)\geqslant0$ 由\cref{lemma:ch08-激波速度符号} 蕴涵 $s=s_i(\tau)\geqslant\lambda_i(W_i(\tau))=\lambda_i(U_+)$，即\eqref{align:ch08-lax-e} 的左半部分。类似论证表明，当 $\tau>0$ 且 $\dot s_i(\cdot)\leqslant0$ 在 $(0,\tau)$ 上时 Lax $E$-条件也成立；但若 $\tau<0$ 且 $\dot s_i(\cdot)<0$ 在 $(\tau,0)$ 上，或 $\tau>0$ 且 $\dot s_i(\cdot)>0$ 在 $(0,\tau)$ 上，则被破坏。上述命题对真非线性情形的推论（此时由\eqref{align:ch08-sdot} 知 $\dot s_i(\cdot)$ 不跨 $0$ 变号）记录于下述

\begin{theorem}[弱激波 Lax 条件的符号]\label{theorem:ch08-弱激波lax}
设 $U_-$ 是系统\eqref{align:ch08-守恒律一维} $i$-特征族的真非线性点，且 $\mathrm{D}\lambda_i(U_-)R_i(U_-)>0$（或 $<0$）。设 $\lambda_i(U_-)$ 是 $\mathrm{D}F(U_-)$ 的单特征值，令 $W_i$ 表示过 $U_-$ 的 $i$-激波曲线，$U_-=W_i(0)$。则一个将 $U_-$ 连接到状态 $U_+=W_i(\tau)$ 的弱 $i$-激波满足 Lax $E$-条件，当且仅当 $\tau<0$（或 $\tau>0$）。
\end{theorem}

于是，在真非线性情形，激波曲线的一半以严格不等式与 Lax $E$-条件\eqref{align:ch08-lax-e} 相容，另一半则不相容。当 $U_-$ 是 $i$-特征场的线性退化点，从而 $\dot s_i(0)=0$ 时，情形更为微妙：若 $\ddot s_i(0)<0$，则 $\dot s_i(\tau)$ 对 $\tau<0$ 为正、对 $\tau>0$ 为负，于是将 $U_-$ 连接到 $U_+=W_i(\tau)$ 的弱 $i$-激波不论 $\tau$ 的符号如何都是容许的。另一方面，若 $\ddot s_i(0)>0$，则 $\dot s_i(\tau)$ 对 $\tau<0$ 为负、对 $\tau>0$ 为正，此时所有（充分）弱 $i$-激波都违反 Lax $E$-条件。如前所述，当 $i$-特征族本身线性退化时，$i$-激波是 $i$-接触间断，以等式满足\eqref{align:ch08-lax-e}。

本节开头启发式分析所暗示的\eqref{align:ch08-压缩}、\eqref{align:ch08-非退化} 与稳定性之间的关系，由下述定理确立：

\begin{theorem}[压缩激波的局部存在性]\label{theorem:ch08-压缩激波存在性}
设系统\eqref{align:ch08-守恒律一维} 严格双曲。考虑初值 $U_0$，使得对 $x\in(-\infty,0)$ 有 $U_0(x)=U_L(x)$，对 $x\in(0,\infty)$ 有 $U_0(x)=U_R(x)$，其中 $U_L$、$U_R$ 是在 $(-\infty,\infty)$ 上与其一阶导数一起有界的光滑函数。进而设左侧状态 $U_-=U_L(0)$ 与右侧状态 $U_+=U_R(0)$ 由一个满足\eqref{align:ch08-压缩} 与\eqref{align:ch08-非退化} 的压缩激波相连。则存在：$T>0$；$[0,T)$ 上的光滑函数 $x=\chi(t)$，$\chi(0)=0$；以及 $(-\infty,\infty)\times[0,T)$ 上的函数 $U$，具有初值 $U_0$ 并满足下述性质。对任一满足 $t\in[0,T)$ 且 $x\neq\chi(t)$ 的 $(x,t)$，$U$ 光滑并在经典意义下满足\eqref{align:ch08-守恒律一维}。进而，对 $t\in[0,T)$，单侧极限 $U(\chi(t)-,t)$ 与 $U(\chi(t)+,t)$ 存在，并由一个速度为 $\dot\chi(t)$ 的压缩激波相连。
\end{theorem}

该证明见\hyperref[section:ch08-注记]{8.8 节}所引文献，它使用沿特征线监控这些函数演化所得到的 $U$ 及其导数的逐点界；即与\cref{theorem:ch07-局部存在与破裂} 的证明同属一类。\eqref{align:ch08-压缩} 与\eqref{align:ch08-非退化} 的作用是确保每个特征族的特征线——或发源于 $x$ 轴、或发源于激波图像——到达上半平面的每一点。可与本节开头的讨论相比较。从极简单的特殊情形 $n=1$ 可获得一些洞见。

于是考虑单个守恒律\eqref{align:ch07-标量守恒律}，并赋予初值 $u_0$，使得 $x\in(-\infty,0)$ 时 $u_0(x)=u_L(x)$，$x\in(0,\infty)$ 时 $u_0(x)=u_R(x)$，其中 $u_L$、$u_R$ 是在 $(-\infty,\infty)$ 上有界且一致 Lipschitz 连续的函数。进而 $u_-=u_L(0)$、$u_+=u_R(0)$ 满足
\begin{align}
f'(u_+)<\frac{f(u_+)-f(u_-)}{u_+-u_-}<f'(u_-).
\label{align:ch08-标量严格lax}
\end{align}
令 $u_-(x,t)$ 与 $u_+(x,t)$ 分别为具有初值 $u_L(x)$、$u_R(x)$ 的\eqref{align:ch07-标量守恒律} 经典解，由\cref{theorem:ch06-6-1-1}，它们在某个 $T>0$ 上定义于 $(-\infty,\infty)\times[0,T)$。在 $[0,T)$ 上，将函数 $\chi$ 定义为常微分方程
\begin{align}
\frac{\mathrm{d}x}{\mathrm{d}t}=\frac{f(u_+(x,t))-f(u_-(x,t))}{u_+(x,t)-u_-(x,t)}
\label{align:ch08-激波运动方程}
\end{align}
满足 $\chi(0)=0$ 的解。最后，在 $(-\infty,\infty)\times[0,T)$ 上定义函数 $u$ 为
\begin{align}
u(x,t)=
\begin{cases}
u_-(x,t),&t\in[0,T),\ x<\chi(t)\\
u_+(x,t),&t\in[0,T),\ x>\chi(t).
\end{cases}
\label{align:ch08-标量解}
\end{align}
显然，对任一满足 $t\in[0,T)$ 且 $x\neq\chi(t)$ 的 $(x,t)$，$u$ 在经典意义下满足\eqref{align:ch07-标量守恒律}。进而 $u(\chi(t)-,t)$ 与 $u(\chi(t)+,t)$ 由速度为 $\dot\chi(t)$ 的激波相连。最后，对充分小的 $T$，Lax $E$-条件
\begin{align}
f'(u(\chi(t)+,t))<\dot\chi(t)<f'(u(\chi(t)-,t)),\qquad t\in[0,T),
\label{align:ch08-标量lax延续}
\end{align}
由连续性成立，因为它在 $t=0$ 处成立。注意，正是由于 Lax $E$-条件，解 $u$ 才仅依赖于初值，即独立于 $x>0$ 时 $u_L(x)$、$x<0$ 时 $u_R(x)$ 所携带的“多余”信息。

如\cref{section:ch07-一维平衡律}所指出，一维守恒律系统既出现在本质上为一维的介质中，也出现在场除一个空间维外均为常数的多空间维介质中。在后一情形下，\cref{theorem:ch08-压缩激波存在性} 确立了平面激波阵面的稳定性，但仅对同样沿法向空间方向变化的扰动而言。自然，研究多空间维平面激波在更广扰动类下的稳定性，以及更一般地研究 $\mathbb R^m$ 中非平面激波阵面的稳定性，是重要的。

\cref{theorem:ch08-压缩激波存在性} 所处理的问题类型，可对 $\mathbb R^m$ 中的守恒律双曲系统\eqref{align:ch04-守恒律组} 表述如下。设 $U_L$、$U_R$ 是 $\mathbb R^m$ 上的光滑函数，$\mathscr F$ 是嵌入 $\mathbb R^m$ 并由其单位法向量场 $v$ 定向的光滑 $(m-1)$ 维超曲面。设 $U_L$ 与 $U_R$ 在 $\mathscr F$ 上的迹 $U_-$、$U_+$ 满足 Rankine--Hugoniot 跳跃条件\eqref{align:ch04-时间差}。记 $U_0$ 为 $\mathbb R^m$ 上在 $\mathscr F$ 负侧与 $U_L$ 一致、在正侧与 $U_R$ 一致的函数。需要构造一个嵌入 $\mathbb R^m\times[0,T)$、在 $t=0$ 处以 $\mathscr F$ 为迹的 $m$ 维超曲面 $\mathscr S$，以及 $\mathbb R^m\times[0,T)$ 上具有初值 $U_0$ 的分片光滑解 $U$，使得对 $(x,t)\notin\mathscr S$ 时 $U$ 光滑。于是 $\mathscr S$ 将是一个由 $\mathscr F$ 演化出的激波。

上述问题已在如下假设下解决。系统\eqref{align:ch04-守恒律组} 配有一致凸熵，并满足某一结构性条件（特别对等熵或非等熵气体动力学 Euler 方程成立）。在 $\mathscr F$ 上每一点（单位法向量为 $v$），\eqref{align:ch08-压缩} 与\eqref{align:ch08-非退化} 以 $\lambda_j(v;U)$、$R_j(v;U)$ 代替 $\lambda_j(U)$、$R_j(U)$ 成立。最后，一组涉及 $U_L$、$U_R$ 的法向导数 $\partial^p U_L/\partial v^p$、$\partial^p U_R/\partial v^p$（最高阶数依赖于 $m$）的复杂相容条件在 $\mathscr F$ 上满足。这些条件是为了避免从 $\mathscr F$ 发出虚假波所必需的。

$\mathscr S$ 与 $U$ 的构造在 Sobolev 空间框架内完成，涉及相当精细的工具（伪微分算子、仿微分演算等）。相关文献列于\hyperref[section:ch08-注记]{8.8 节}。

另一个值得关注的严重问题是激波的内部稳定性。结果表明，只要系统真非线性且激波为弱激波，Lax $E$-条件在该方向上也是有效的；然而在更一般的情形下它并不充分。为此，我们必须考虑附加的、更具选择性的激波容许性准则，它们将在以下各节中引入。

\end{document}
